\subsection{Three small affine obstructions for arbitrary acts}

\begin{lemma}[A local contracted algebra with two columns]
\label{C:local-two-column-act-cups}
Let $M$ have a completely simple minimal ideal with two
columns, and suppose its contracted rational algebra
$A=\mathbb QM/\mathbb QK$ is local.  Every rational degree-two
product of every finite act cone of $M$ is zero.
There is no bound on the number of rows or states, and the
minimal maximal group may be nontrivial.
\end{lemma}
\begin{proof}
Every distinguished nonzero basis element of $A$ has a high
power that is either zero or a distinguished idempotent basis
element.  A local algebra has no idempotents other than zero
and one.  If that power is one, the original basis element is
an actual finite unit in the pointed quotient monoid $M/K$.
A nontrivial cyclic subgroup would give the idempotent average
of its distinct basis elements, which is neither zero nor one.
This is impossible.  Consequently every outside basis element
other than the identity is nilpotent.

The old column augmentation defines the character
$\chi(m)=\tau_m(1)-\tau_m(0)$ of $A$.
It is zero on every such nilpotent basis element.
All nonidentity actors therefore have constant column maps.
The identity is the only actor with nonconstant columns, so
any row is a common fixed row for those actors.
Proposition~\ref{C:act-common-affine-origin-cups} applies.
\end{proof}

\begin{lemma}[Two quotient covectors on at most five rows]
\label{C:two-covector-affine-origin}
Let $M$ have a minimal ideal with at most five rows and two
columns.  Put $V=\widetilde{\mathbb QI}$ and let $\chi$ be its
column augmentation character.  If
\[
 \dim_{\mathbb Q}\operatorname{Hom}_A(V,\mathbb Q_\chi)\geq2,
 \qquad A=\mathbb QM/\mathbb QK,
\]
then every rational degree-two product of every finite act
cone is zero.  No bounded-homology assumption is needed.
\end{lemma}
\begin{proof}
Choose two independent quotient covectors and evaluate them
on $[i]-[i_0]$.  This gives a map $f:I\to\mathbb Q^2$
whose image $F$ affinely spans the plane.  Each row actor has
\[
 f(\sigma_m i)=\chi(m)f(i)+b_m.
\]
For $\chi(m)=1$, the translation $b_m$ preserves the finite
set $F$ and hence is zero.  Such an actor preserves every
fiber of $f$.  Every actor with $\chi(m)=-1$ induces a central
reflection of $F$.  All these reflections are the same:
the composite of two is a translation of $F$, necessarily
the identity.

If there is no negative actor, then $|F|\geq3$ and $|I|\leq5$
give a singleton fiber.  Its row is fixed by every actor with
nonconstant columns.  If there is a negative actor, the
centrally symmetric set $F$ has at least four points, because
it spans a plane.  With at most five rows, at most one fiber
can contain more than one point.  Thus there is an opposite
pair of singleton fibers, or a singleton central fiber.
In the first case the midpoint of the two corresponding row
basis vectors is fixed by every nonconstant-column actor;
in the second case the corresponding row is fixed.
Either gives the common rational affine origin required by
Proposition~\ref{C:act-common-affine-origin-cups}.
\end{proof}

\begin{lemma}[Three outside elements, at every row and state count]
\label{C:three-outside-all-two-column-act-cups}
If a finite monoid has a minimal ideal with two columns and
at most three outside elements, all rational degree-two
products of all its finite act cones vanish.
\end{lemma}
\begin{proof}
Let $P$ be the actors with nonconstant column map.  This is an
actual submonoid containing the identity: its column slopes
multiply in $\{1,-1\}$, so a product remains outside $K$.
Its order is at most three.

Every row action of such a monoid has a common rational affine
fixed point except possibly when $P$ consists of the identity
adjoined to a two-element left-zero semigroup.  Here is a direct
verification.  If $P$ is a group, average a row orbit.  If its
minimal ideal has one element, that zero's row image is fixed
pointwise by $P$.  A two-element minimal ideal is either a
group of order two, a left-zero band, or a right-zero band.
For the group use its local group average; the only extra
element is the global identity.  For the right-zero band,
$\sigma_x\sigma_e=\sigma_e$ makes every point of
$\operatorname{im}\sigma_e$ fixed.  These cases also cover
orders one and two.

Only the two-left-zero case remains.  It uses all three
outside elements, and all their products stay outside $K$.
Consequently $A=\mathbb QP$, the unitalization of the algebra
of two left zeros.  If these are $p_0,p_1$, the orthogonal
idempotents $p_0,1-p_0$ and radical element $p_1-p_0$ identify
$A$ with the three-dimensional path algebra of one arrow.
It is hereditary, so every row module has projective dimension
at most one.  Lemma~\ref{C:row-pd-one-act-cups} applies in
this last case, and the affine-origin proposition applies in
all other cases.
\end{proof}
